%==============================================
% VOLUME III: ARRAY ODYSSEY
% (was: "Part 3")
%==============================================
%
% NOTE ON VOLUME STRUCTURE:
%   Arliz is now organized as three volumes rather than three parts of
%   a single book (see Preface and Introduction). This file represents
%   the content of Volume III, the destination of the whole work. It
%   assumes the representational vocabulary of Volume I (BIT GENESIS) and the hardware model of Volume II (Architecture &
%   Circuitry). If compiled as a standalone book, replace \part{} below
%   with a suitable top-level structure for its own document.
%
% MERGE NOTES (step by step, voltage -> bits -> arrays):
%   - Keeps ALL original Part 3 chapters (theory, memory layout,
%     traversal, dynamic arrays, multidimensional arrays, sparse arrays,
%     bit arrays, searching, sorting, rotation/permutation, two-pointer,
%     prefix sums, views/slices, broadcasting, scan/reduce, etc.)
%   - DROPS the memory/storage hardware-testing and diagnostic chapters
%     that were appended to the end of the original Part 3 (March
%     algorithms for RAM testing, SMART attributes, etc.) -- these belong
%     to hardware repair, not array theory, and have no home in this work.
%   - MERGES IN from old Part 4 (Data Structures & Algorithms), placed
%     immediately after the core array-odyssey material, since every one
%     of these structures and algorithms IS an array technique:
%       * Stacks/Queues/Deques/Priority Queues/Heaps
%       * Hash tables and probabilistic structures (Bloom filters, etc.)
%       * Strings, pattern matching, suffix structures, tries
%       * Graph representations and classic graph algorithms
%       * Range-query structures (segment tree, Fenwick tree, sparse table)
%       * Monotonic stack/queue
%       * Dynamic programming
%       * Divide & conquer, greedy, backtracking
%       * Online/streaming algorithms, amortized analysis,
%         cache-oblivious algorithms
%   - MERGES IN from old Part 5 (Parallelism & Systems), the
%     algorithm-level (not pure-hardware) parallel chapters, placed as
%     a "parallel processing of arrays" section near the end:
%       * Concurrency fundamentals, performance laws (Amdahl/Gustafson)
%       * Threads, thread pools, work stealing, fork-join
%       * OpenMP
%       * Synchronization primitives, lock-free structures
%       * GPU array algorithms: sorting, reduction, scan
%       * Distributed array processing: MPI, partitioning, MapReduce,
%         Spark, PGAS, consensus
%       * Real-time and energy-aware array processing
%   - MERGES IN from old Part 6 (Synthesis & Frontiers), the
%     application-domain array chapters, placed as the final "where
%     arrays are used today" section:
%       * Tensors, ML array operations, attention/transformers
%       * Linear algebra (BLAS/LAPACK), scientific computing, FFT
%       * Image/audio/video/signal processing
%       * Bioinformatics array applications
%       * Quantum state vectors
%       * Array programming languages and ecosystems (APL, NumPy, etc.)
%       * Persistent/compressed/succinct array representations
%       * Privacy-preserving and distributed-ledger array structures
%       * Emerging array hardware/paradigms (neuromorphic, optical, DNA)
%
\part{Array Odyssey}
\label{part:array-odyssey}

\begin{partintro}
\lettrine[lines=3]{A}{t last} we arrive at arrays themselves---the most fundamental data structure in computing, and the destination of everything Volumes I and II built up to. This volume explores arrays from every conceivable angle: theory, implementation, optimization, the data structures and algorithms built on top of them, the parallel and distributed systems that process them at scale, and the application domains---from machine learning to quantum computing---where arrays do their real work.

\vspace{1em}
\textbf{What Makes This Different:}
\begin{itemize}[noitemsep]
    \item \textbf{Complete Coverage:} Every variant, every technique, every structure built on arrays
    \item \textbf{Hardware Awareness:} Cache effects, alignment, SIMD, and GPU execution---grounded in Volume II
    \item \textbf{Rigorous Analysis:} Mathematical foundations and performance, from a single access to a distributed cluster
    \item \textbf{Practical Optimization:} Making arrays fast in real systems, from one core to many
\end{itemize}

\begin{quote}
\textit{``Simplicity is prerequisite for reliability.''}

\hfill--- \textsc{Edsger W. Dijkstra}
\end{quote}
\end{partintro}

% ============================================================
% SECTION A: Array theory and memory layout (original Part 3, kept as-is)
% ============================================================

\chapter{Mathematical Foundations of Arrays}
Formal definition, array as function from index set to value set, properties, mathematical operations.

\chapter{Array as Abstract Data Type}
Interface specification, invariants, preconditions, postconditions, ADT vs. implementation.

\chapter{One-Dimensional Array: Memory Layout}
Contiguous allocation, base address, element offset calculation, memory alignment.

\chapter{Address Calculation and Pointer Arithmetic}
Computing element addresses, stride, pointer arithmetic rules, bounds.

\chapter{Zero-Based vs. One-Based Indexing}
Historical reasons, language conventions, conversion, pros and cons.

\chapter{Array Bounds and Boundary Conditions}
Index range, out-of-bounds access, undefined behavior, bounds checking strategies.

\chapter{Array Allocation: Stack vs. Heap}
Automatic (stack) allocation, dynamic (heap) allocation, lifetime management, size limitations.

\chapter{Static Array Declaration}
Compile-time size, storage duration, initialization, usage patterns.

\chapter{Dynamic Memory Allocation for Arrays}
malloc/calloc/realloc, new/delete, memory leaks, fragmentation.

\chapter{Array Deallocation and Resource Management}
Proper cleanup, RAII pattern, smart pointers, memory ownership.

\chapter{Array Initialization Techniques}
Zero-initialization, value initialization, aggregate initialization, designated initializers.

\chapter{Array Initialization Performance}
memset internals, compiler optimizations, initialization costs, lazy initialization.

\chapter{Array Copying Fundamentals}
Shallow vs. deep copy, copy semantics, assignment operators.

\chapter{Memory Transfer Operations}
memcpy, memmove, bulk copy performance, vectorized copying.

\chapter{Array Access Patterns and Performance}
Sequential access, random access, strided access, impact on cache and prefetching.

\chapter{Cache Line Utilization}
Cache line size (64 bytes typical), spatial locality, utilizing full cache lines.

\chapter{Prefetching for Array Access}
Hardware prefetching, software prefetching, prefetch instructions, prefetch distance.

\chapter{Array Traversal Optimization}
Loop order, blocking/tiling, minimizing cache misses, memory access patterns.

\chapter{Dynamic Arrays: Implementation}
Capacity vs. size, growth strategy, amortized analysis, C++ std::vector internals.

\chapter{Dynamic Array Growth Strategies}
Doubling, multiplicative (golden ratio), additive growth, tradeoff analysis.

\chapter{Dynamic Array Amortized Analysis}
Proving O(1) amortized append, potential method, aggregate analysis.

\chapter{Dynamic Array Reallocation}
In-place vs. move to new location, realloc behavior, minimizing copying.

\chapter{Dynamic Array Shrinking}
Shrink-to-fit, hysteresis, when to shrink, avoiding oscillation.

\chapter{Dynamic Arrays in Standard Libraries}
C++ std::vector, Java ArrayList, Python list, implementation comparison.

\chapter{Dynamic Array: Iterator Invalidation}
When iterators become invalid, reallocation impact, safe iteration patterns.

\chapter{Multidimensional Array Fundamentals}
Dimensions, shape, indexing multiple dimensions, row-major vs. column-major.

\chapter{Row-Major Order (C Convention)}
Layout in memory, linearization formula, cache implications.

\chapter{Column-Major Order (Fortran Convention)}
Layout differences, when it matters, interfacing with Fortran libraries (BLAS, LAPACK).

\chapter{Multidimensional Array Address Calculation}
Linear index from multiple indices, stride computation, optimization.

\chapter{Multidimensional Array Traversal}
Loop nesting order, cache-friendly traversal, impact of memory layout.

\chapter{Array of Arrays vs. True Multidimensional}
Pointer indirection, flexibility, memory overhead, cache implications.

\chapter{Jagged Arrays}
Non-rectangular arrays, array of arrays for irregular sizes, applications.

\chapter{Dope Vectors and Array Descriptors}
Metadata representation, bounds information, enabling flexible array operations.

\chapter{Sparse Array Fundamentals}
When most elements are zero (or default), compression ratio, sparsity pattern.

\chapter{Sparse Array: Coordinate List (COO)}
List of (row, col, value) tuples, easy construction, inefficient access.

\chapter{Sparse Array: Compressed Sparse Row (CSR)}
Row pointers, column indices, values, efficient row operations.

\chapter{Sparse Array: Compressed Sparse Column (CSC)}
Column-oriented version of CSR, efficient column operations.

\chapter{Sparse Array: Dictionary of Keys (DOK)}
Hash table representation, flexible construction, conversion to other formats.

\chapter{Sparse Array: List of Lists (LIL)}
List per row, incremental construction, conversion to CSR.

\chapter{Sparse Matrix Operations}
Sparse matrix-vector multiply, sparse-sparse operations, fill-in problem.

\chapter{Bit Arrays: Representation}
Packing bits into words, storage efficiency, applications (sets, flags).

\chapter{Bit Array: Indexing and Access}
Computing word and bit offset, extracting bits, setting bits.

\chapter{Bit Array Operations}
Bitwise operations on arrays, parallel operations, applications.

\chapter{Bitmap Indexes}
Using bitmaps for indexing, rank and select operations, succinct data structures preview.

\chapter{Circular Arrays and Ring Buffers}
Wrap-around indexing, modulo arithmetic, avoiding data movement, queue implementation.

\chapter{Circular Array: Head and Tail Management}
Tracking fill level, full vs. empty distinction, producer-consumer patterns.

% ============================================================
% SECTION B: Searching and sorting (original Part 3, kept as-is)
% ============================================================

\chapter{Array Searching: Linear Search}
Sequential scan, early termination, best/worst/average case, sentinel technique.

\chapter{Array Searching: Binary Search}
Prerequisites (sorted array), algorithm, invariants, complexity analysis.

\chapter{Binary Search: Implementation Variants}
Iterative vs. recursive, overflow-safe midpoint, lower\_bound vs. upper\_bound.

\chapter{Binary Search: Applications}
Finding insertion point, counting occurrences, rotated arrays.

\chapter{Interpolation Search}
Estimating position, prerequisites (uniform distribution), O(log log n) average case.

\chapter{Exponential Search}
Unbounded search, doubling range, combination with binary search.

\chapter{Ternary Search}
Dividing into three parts, unimodal functions, finding extrema.

\chapter{Jump Search}
Block jumping, optimal jump size, combination with linear search, cache-friendly properties.

\chapter{Fibonacci Search}
Division using Fibonacci numbers, single comparison per step, applications.

\chapter{Sorting Fundamentals}
Comparison-based lower bound, stability, adaptivity, in-place vs. out-of-place.

\chapter{Bubble Sort}
Adjacent swaps, optimization (early termination), adaptive behavior, teaching value.

\chapter{Selection Sort}
Finding minimum, swap to position, in-place but not stable, O(n\textsuperscript{2}) always.

\chapter{Insertion Sort}
Inserting into sorted prefix, adaptive (O(n) on nearly sorted), online algorithm.

\chapter{Shell Sort}
Gap sequences, h-sorting, diminishing increments, performance depends on gap sequence.

\chapter{Merge Sort Fundamentals}
Divide and conquer, merge procedure, guaranteed O(n log n), stability.

\chapter{Merge Sort: Implementation Variants}
Top-down recursive, bottom-up iterative, in-place merging attempts, natural merge sort.

\chapter{Merge Sort: Optimization Techniques}
Switching to insertion sort for small subarrays, eliminating copy to auxiliary array.

\chapter{External Merge Sort}
Sorting data larger than memory, multi-way merge, minimizing I/O, B-tree connection.

\chapter{Quick Sort Fundamentals}
Partitioning, divide and conquer, average O(n log n), worst O(n\textsuperscript{2}).

\chapter{Quick Sort: Pivot Selection Strategies}
First element, random element, median-of-three, median-of-medians, ninther.

\chapter{Quick Sort: Partitioning Schemes}
Lomuto partition, Hoare partition, three-way partition (Dutch National Flag).

\chapter{Quick Sort: Optimizations}
Tail recursion elimination, switching to insertion sort, handling duplicates.

\chapter{Randomized Quick Sort}
Random pivot selection, probabilistic analysis, expected O(n log n).

\chapter{Introsort}
Hybrid algorithm, depth limiting, switching to heapsort, C++ std::sort implementation.

\chapter{Heap Sort Fundamentals}
Building heap, extract-max repeatedly, in-place, not stable, guaranteed O(n log n).

\chapter{Binary Heap: Array Representation}
Complete binary tree in array, parent/child index formulas, level-order storage.

\chapter{Heapify Operations}
Sift-up (bubble-up), sift-down (percolate-down), building heap in O(n).

\chapter{Heap Sort: Optimization}
Bottom-up heap construction, cache optimization, comparison count.

\chapter{Counting Sort}
Integer sorting, counting occurrences, cumulative sum, linear time O(n+k).

\chapter{Radix Sort Fundamentals}
Digit-by-digit sorting, LSD vs. MSD, stable sub-sorting required.

\chapter{Radix Sort: Implementations}
Base/radix selection, counting sort as sub-routine, in-place variants.

\chapter{Bucket Sort}
Distributing into buckets, uniform distribution assumption, O(n) average case.

\chapter{Timsort}
Python's sorting algorithm, galloping mode, merge runs, adaptive merging.

\chapter{pdqsort (Pattern-Defeating Quicksort)}
Rust's sort, detecting bad patterns, fallback strategies, block partitioning.

\chapter{Sorting Networks}
Comparator networks, Batcher's sort, bitonic sort, parallel sorting.

% ============================================================
% SECTION C: In-place manipulation, two pointers, prefix sums
%   (original Part 3, kept as-is)
% ============================================================

\chapter{Array Rotation: Reversal Algorithm}
Reverse portions, three reversals, in-place, O(n) time, elegant proof.

\chapter{Array Rotation: Juggling Algorithm}
GCD-based, moving elements in cycles, in-place, optimal moves.

\chapter{Array Rotation: Block Swap Algorithm}
Recursive block swapping, handling unequal blocks, in-place rotation.

\chapter{Array Reversal In-Place}
Two-pointer technique, swapping elements, O(n) time, O(1) space.

\chapter{Array Permutation: Application}
Applying permutation in-place, cycle decomposition, O(n) time.

\chapter{Array Permutation: Generation}
Generating all permutations, lexicographic order, Heap's algorithm, iterative methods.

\chapter{Array Shuffling}
Fisher-Yates shuffle, uniform random permutation, in-place, bias-free.

\chapter{Array Partitioning Algorithms}
Two-way partition (quicksort), three-way partition (Dutch flag), k-way partition.

\chapter{Stable Partitioning}
Maintaining relative order, applications, implementation strategies.

\chapter{Quickselect Algorithm}
Finding k-th smallest element, partition-based selection, average O(n).

\chapter{Median of Medians}
Deterministic linear-time selection, choosing good pivots, theoretical importance.

\chapter{Two-Pointer Technique: Basics}
Meeting in middle, opposite directions, applications (two sum, container problems).

\chapter{Two-Pointer: Sliding Window Fixed Size}
Window size k, maintaining window state, deque for min/max in window.

\chapter{Two-Pointer: Sliding Window Variable Size}
Expanding and contracting, maintaining invariants, applications (substring problems).

\chapter{Three-Pointer and Multiple-Pointer Techniques}
Dutch National Flag, merging sorted arrays, complex partitioning.

\chapter{Prefix Sum: Construction}
Cumulative sum array, O(n) construction, range sum query in O(1).

\chapter{Prefix Sum: Applications}
Range queries, subarray sum, equilibrium index, difference arrays.

\chapter{Two-Dimensional Prefix Sum}
2D cumulative sum, O(1) rectangle sum query, inclusion-exclusion.

\chapter{Difference Arrays}
Inverse of prefix sum, range updates in O(1), final array in O(n).

\chapter{Kadane's Algorithm}
Maximum subarray sum, dynamic programming, linear time, variations.

\chapter{Kadane's Algorithm: Variations}
Maximum product subarray, circular array, maximum sum with deletions.

\chapter{Subarray Problems}
Subarray vs. subsequence, contiguous vs. non-contiguous, counting subarrays.

\chapter{Dutch National Flag Problem}
Three-way partitioning, single pass, constant space, applications to quicksort.

\chapter{Array Intersection}
Finding common elements, sorted arrays (two-pointer), unsorted (hash set).

\chapter{Array Union and Set Operations}
Union, intersection, difference, symmetric difference, complexity analysis.

\chapter{Merging Sorted Arrays}
Two-way merge, k-way merge (heap-based), in-place merging.

\chapter{Array Compaction}
Removing elements while maintaining order, in-place, stable compaction.

\chapter{Array Gap Removal}
Removing duplicates, removing specific elements, in-place techniques.

\chapter{Array Element Frequency}
Counting occurrences, mode finding, frequency map, Boyer-Moore majority vote.

\chapter{Majority Element Algorithm}
Boyer-Moore voting, linear time, constant space, verification step.

\chapter{Array Rearrangement Problems}
Alternating positive-negative, even-odd segregation, custom ordering.

\chapter{In-Place Array Manipulation Techniques}
Swapping, reversing, rotating without extra space, index mapping.

\chapter{Constant Space Algorithms}
O(1) auxiliary space, in-place transformations, iterative approaches.

% ============================================================
% SECTION D: Alignment, bounds checking, views, broadcasting
%   (original Part 3, kept as-is)
% ============================================================

\chapter{Array Memory Alignment: Cache Lines}
Aligning to 64-byte boundaries, false sharing prevention, performance impact.

\chapter{Array Memory Alignment: SIMD}
16-byte, 32-byte, 64-byte alignment requirements, aligned\_alloc, posix\_memalign.

\chapter{Alignment Attributes and Directives}
Compiler attributes (alignas, \_\_attribute\_\_), pragma directives, ensuring alignment.

\chapter{Array Bounds Checking Strategies}
Runtime bounds checking, compiler-generated checks, safe array access.

\chapter{Bounds Checking: Performance Impact}
Overhead analysis, optimization opportunities, eliminating redundant checks.

\chapter{Bounds Checking: Language Support}
Rust's borrow checker, C++ std::array, Java array bounds, safe languages.

\chapter{Array Views and Slices}
Non-owning references, NumPy slicing, stride manipulation, views vs. copies.

\chapter{Array Slicing Semantics}
Syntax (start:stop:step), negative indices, omitted bounds, Python/NumPy conventions.

\chapter{Strided Array Access}
Non-contiguous access, stride specification, performance implications, vectorization challenges.

\chapter{Array Subviews and Windows}
Creating windows without copying, sliding windows, overlapping windows.

\chapter{Array Broadcasting}
Implicit dimension expansion, NumPy broadcasting rules, shape compatibility.

\chapter{Broadcasting: Implementation}
Virtual expansion, avoiding memory duplication, optimized iteration.

\chapter{Array Reduction Operations}
Sum, product, min, max, logical operations, custom reductions.

\chapter{Array Scan Operations}
Inclusive scan (cumulative sum), exclusive scan (shifted cumulative sum).

\chapter{Array Gather Operations}
Indexed reads, gather from array using index array, vectorization support.

\chapter{Array Scatter Operations}
Indexed writes, scatter to array using index array, conflict handling.

\chapter{Gather-Scatter in SIMD}
AVX2 gather instructions, AVX-512 gather/scatter, performance characteristics.

\chapter{Array Comprehensions}
Declarative array construction, Python list comprehensions, functional style.

\chapter{Array Generators and Iterators}
Lazy evaluation, generator expressions, memory efficiency, iterator protocols.

\chapter{Array Reshaping}
Changing dimensions without copying data, view manipulation, Fortran vs. C order.

\chapter{Array Transposition}
Swapping dimensions, in-place transpose (square), out-of-place (rectangular).

\chapter{Cache-Oblivious Transpose}
Recursive divide-and-conquer, optimal cache usage, no tuning parameters.

\chapter{Array Concatenation}
Joining arrays along axis, horizontal stack (column-wise), vertical stack (row-wise).

\chapter{Array Splitting}
Dividing array, chunk-based splitting, equal vs. unequal splits.

\chapter{Array Tiling}
Dividing into tiles, blocked algorithms, cache optimization, matrix multiplication.

\chapter{Array Padding Techniques}
Border padding (image processing), alignment padding, ghost cells (stencil computations).

\chapter{Toeplitz and Circulant Matrices}
Special array structures, efficient storage, fast multiplication (FFT).

\chapter{Triangular and Symmetric Matrices}
Storing only half, index mapping, packed storage, BLAS conventions.

\chapter{Band Matrices and Sparse Diagonal}
Diagonal storage format, bandwidth, tridiagonal systems, efficient solvers.

\chapter{Array Vectorization Fundamentals}
Writing vector-friendly code, enabling autovectorization, compiler hints.

\chapter{Loop Vectorization}
Data dependencies, trip count, alignment, remainder loops.

\chapter{Array Expressions and Lazy Evaluation}
Expression templates (C++), deferred computation, Eigen library example.

\chapter{Array Versioning and Copy-on-Write}
Persistent arrays, sharing data between versions, efficient versioning.

\chapter{Immutable Arrays}
Functional programming, structural sharing, immutability benefits.

\chapter{Array Performance Profiling}
Measuring access patterns, cache simulation tools, hardware counters.

\chapter{Roofline Model for Arrays}
Arithmetic intensity, bandwidth bounds, computational bounds, optimization guidance.

\chapter{Array Benchmarking Methodology}
Microbenchmarks, cache warming, timing techniques, statistical significance.

% ============================================================
% SECTION E: Foundational structures built on arrays
%   (merged from old Part 4)
% ============================================================

\chapter{Abstract Data Types}
ADT specification, interface vs. implementation, invariants, contracts.

\chapter{Stacks: Array Implementation}
LIFO principle, push/pop operations, array-based stack, dynamic resizing, applications.

\chapter{Stack Applications}
Expression evaluation, parenthesis matching, backtracking, function call stack.

\chapter{Monotonic Stack}
Maintaining monotonic property, next greater element, O(n) amortized.

\chapter{Monotonic Stack Applications}
Largest rectangle in histogram, stock span problem, maximum area problems.

\chapter{Queues: Circular Array Implementation}
FIFO principle, circular buffer, enqueue/dequeue, front/rear pointers, full vs. empty detection.

\chapter{Queue Applications}
BFS, task scheduling, buffering, producer-consumer patterns.

\chapter{Deques: Double-Ended Queues}
Operations at both ends, circular array implementation, applications (sliding window).

\chapter{Monotonic Queue}
Sliding window maximum/minimum, deque-based, O(n) total operations.

\chapter{Priority Queue: Concept}
Priority-based access, operations (insert, find-min, delete-min), applications.

\chapter{Binary Heaps: Array Representation}
Complete binary tree in array, parent-child formulas, heap property (min-heap, max-heap).

\chapter{Binary Heap: Operations}
Insert (sift-up), extract-min (sift-down), decrease-key, build-heap in O(n).

\chapter{Binary Heap: Analysis}
Time complexity analysis, comparison count, space efficiency.

\chapter{Heap Variants}
d-ary heaps, binary heaps vs. d-ary heaps, cache performance.

\chapter{Heap Sort Revisited}
Using heap as sorting technique, in-place sorting, comparison with other sorts.

\chapter{Priority Queue Applications}
Dijkstra's algorithm, Prim's MST, event simulation, job scheduling.

% --- Hash tables and probabilistic structures ---

\chapter{Hash Table Fundamentals}
Hash function, collision resolution, load factor, dynamic resizing.

\chapter{Hash Functions}
Division method, multiplication method, universal hashing, perfect hashing.

\chapter{Collision Resolution: Chaining}
Separate chaining with linked lists, array of lists, dynamic arrays for chains.

\chapter{Collision Resolution: Open Addressing}
Linear probing, quadratic probing, double hashing, clustering problems.

\chapter{Open Addressing: Advanced Techniques}
Robin Hood hashing, hopscotch hashing, cuckoo hashing (multiple tables).

\chapter{Hash Table Resizing}
Doubling table size, rehashing all elements, amortized analysis, incremental resizing.

\chapter{Hash Table Performance Analysis}
Load factor impact, expected chain length, probe sequence length, worst-case scenarios.

\chapter{Perfect Hashing}
Collision-free for static sets, two-level hashing, construction algorithms.

\chapter{Minimal Perfect Hashing}
Space-optimal perfect hashing, CHD algorithm, MPHF construction.

\chapter{Bloom Filters: Fundamentals}
Probabilistic set membership, bit array with k hash functions, false positive rate.

\chapter{Bloom Filters: Analysis}
Optimal k (number of hash functions), size determination, false positive probability.

\chapter{Bloom Filter Variants}
Counting Bloom filters, scalable Bloom filters, cuckoo filters.

\chapter{Count-Min Sketch}
Frequency estimation, point queries, range queries, space-accuracy tradeoffs.

\chapter{HyperLogLog}
Cardinality estimation, approximate distinct counting, logarithmic space.

\chapter{Approximate Membership and Counting}
Space-efficient probabilistic structures, applications (databases, networks, caches).

% --- Strings, pattern matching, and text-array structures ---

\chapter{Strings as Character Arrays}
Null termination (C strings), length-prefixed strings, string representations.

\chapter{String Manipulation Primitives}
Concatenation, substring, character access, modification operations.

\chapter{String Matching: Naive Algorithm}
Brute force, sliding window, complexity O(nm), when it's acceptable.

\chapter{Knuth-Morris-Pratt (KMP) Algorithm}
Failure function, avoiding redundant comparisons, O(n+m) complexity, preprocessing.

\chapter{Boyer-Moore Algorithm}
Bad character rule, good suffix rule, right-to-left scanning, sublinear average case.

\chapter{Boyer-Moore-Horspool Variant}
Simplified bad character rule, practical performance, implementation simplicity.

\chapter{Rabin-Karp Algorithm}
Rolling hash, fingerprint matching, multiple pattern matching, collision handling.

\chapter{Aho-Corasick Algorithm}
Multiple pattern matching, trie with failure links, dictionary matching applications.

\chapter{Suffix Arrays: Construction}
Sorting suffixes, O(n log n) construction algorithms, linear-time construction (SA-IS).

\chapter{Suffix Arrays: Applications}
Pattern matching, longest common substring, substring queries.

\chapter{LCP Array (Longest Common Prefix)}
Auxiliary array for suffix array, Kasai's algorithm, applications to string problems.

\chapter{Suffix Trees vs. Suffix Arrays}
Space-time tradeoffs, construction complexity, practical considerations.

\chapter{Tries: Array-Based Nodes}
Alphabet array in each node, space overhead, fast lookup.

\chapter{Compressed Tries and Radix Trees}
Path compression, patricia trees, space efficiency.

\chapter{Trie Applications}
Autocomplete, spell checking, IP routing, prefix matching.

\chapter{Ternary Search Trees}
Space-efficient trie variant, BST-like structure, balanced performance.

% --- Graphs on arrays ---

\chapter{Graph Representations: Overview}
Adjacency matrix, adjacency list, edge list, tradeoffs for different operations.

\chapter{Adjacency Matrix Representation}
2D array, space O(V\textsuperscript{2}), edge checking O(1), suitable for dense graphs.

\chapter{Matrix Operations on Graphs}
Transitive closure (Floyd-Warshall), matrix multiplication for path counting.

\chapter{Adjacency List: Array of Lists}
Space O(V+E), neighbor iteration, suitable for sparse graphs.

\chapter{Adjacency List: Implementation Variants}
Array of vectors, array of linked lists, array of dynamic arrays.

\chapter{Edge List Representation}
Array of edges, simple representation, useful for edge-centric algorithms.

\chapter{Graph Algorithm Fundamentals}
Traversal, path finding, connectivity, topological sort, minimum spanning tree.

\chapter{Breadth-First Search (BFS) on Arrays}
Queue-based traversal, level-order exploration, shortest path in unweighted graphs.

\chapter{Depth-First Search (DFS) on Arrays}
Stack-based (explicit or recursive), preorder/postorder, cycle detection.

\chapter{Topological Sorting}
DFS-based, Kahn's algorithm (BFS-based), applications (dependency resolution).

\chapter{Shortest Path: Dijkstra's Algorithm}
Priority queue (binary heap), relaxation, non-negative weights, O((V+E) log V).

\chapter{Shortest Path: Bellman-Ford Algorithm}
Handling negative weights, negative cycle detection, O(VE) complexity.

\chapter{All-Pairs Shortest Paths: Floyd-Warshall}
Dynamic programming on adjacency matrix, O(V\textsuperscript{3}), dense graph suitability.

\chapter{Minimum Spanning Tree: Prim's Algorithm}
Priority queue, growing MST from single vertex, O((V+E) log V).

\chapter{Minimum Spanning Tree: Kruskal's Algorithm}
Sorting edges, union-find for cycle detection, O(E log E).

\chapter{Disjoint-Set Union (Union-Find)}
Array-based parent pointers, path compression, union by rank/size.

\chapter{Union-Find: Analysis}
Inverse Ackermann function, nearly constant amortized time, proof sketch.

\chapter{Union-Find: Applications}
Connected components, Kruskal's MST, dynamic connectivity, percolation.

% --- Range-query structures ---

\chapter{Segment Trees: Fundamentals}
Binary tree for range queries, array representation, recursive structure.

\chapter{Segment Tree: Construction and Queries}
Building O(n), range query O(log n), point update O(log n).

\chapter{Segment Tree: Lazy Propagation}
Range updates, deferred updates, maintaining correctness, complexity analysis.

\chapter{Segment Tree: Variations}
Range minimum query (RMQ), range sum, range GCD, arbitrary associative operations.

\chapter{Fenwick Tree (Binary Indexed Tree)}
Cumulative frequency, update and query O(log n), compact representation.

\chapter{Fenwick Tree: Implementation}
Bit manipulation for parent/child, 1-indexed arrays, prefix sum queries.

\chapter{Fenwick Tree: Range Updates}
Difference array technique, point query with range update, dual representation.

\chapter{Sparse Table}
Preprocessing for RMQ, O(n log n) space, O(1) query, immutable arrays.

\chapter{Sparse Table: Construction}
Dynamic programming, overlapping ranges, idempotent operation requirement.

\chapter{Range Query Structures Comparison}
Segment tree, Fenwick tree, sparse table, tradeoffs and use cases.

% --- Algorithmic design paradigms on arrays ---

\chapter{Dynamic Programming Fundamentals}
Optimal substructure, overlapping subproblems, memoization vs. tabulation.

\chapter{DP: State Representation in Arrays}
1D, 2D, multi-dimensional state spaces, state transition.

\chapter{DP: Space Optimization}
Rolling arrays, reducing dimensions, in-place updates.

\chapter{0/1 Knapsack Problem}
Items with weight and value, capacity constraint, 2D DP array.

\chapter{Unbounded Knapsack}
Unlimited copies of items, 1D DP array suffices, coin change connection.

\chapter{Subset Sum Problem}
Boolean DP, target sum reachability, space-optimized solution.

\chapter{Longest Common Subsequence (LCS)}
2D DP table, traceback for reconstruction, O(nm) time and space.

\chapter{Longest Increasing Subsequence (LIS)}
O(n\textsuperscript{2}) DP solution, O(n log n) binary search solution, patience sorting.

\chapter{Edit Distance (Levenshtein)}
2D DP, insertions/deletions/substitutions, sequence alignment.

\chapter{Edit Distance: Space Optimization}
Linear space, Hirschberg's algorithm, divide and conquer.

\chapter{Matrix Chain Multiplication}
Optimal parenthesization, 2D DP, O(n\textsuperscript{3}) complexity, reconstruction.

\chapter{Optimal Binary Search Trees}
Weighted search cost, 2D DP, Knuth's optimization.

\chapter{Palindrome Problems}
Longest palindromic substring, palindrome partitioning, DP on strings.

\chapter{DP on Grids and Paths}
Minimum path sum, unique paths, grid DP patterns.

\chapter{DP with Bitmasks}
Representing subsets, traveling salesman, Hamiltonian path, assignment problem.

\chapter{Divide and Conquer Fundamentals}
Recurrence relations, master theorem, merge step complexity.

\chapter{Merge Sort Analysis}
Recurrence T(n) = 2T(n/2) + O(n), master theorem application.

\chapter{Quicksort Analysis}
Average case, randomized analysis, worst case handling.

\chapter{Binary Search Analysis}
Recurrence T(n) = T(n/2) + O(1), logarithmic complexity.

\chapter{Closest Pair of Points}
Divide and conquer, O(n log n), geometric algorithms on arrays.

\chapter{Convex Hull Algorithms}
Graham scan, Jarvis march, quick hull, geometric array processing.

\chapter{Computational Geometry Primitives}
Point arrays, line segments, orientation test, intersection tests.

\chapter{Greedy Algorithms on Arrays}
Greedy choice property, optimal substructure, exchange arguments.

\chapter{Activity Selection Problem}
Interval scheduling, earliest finish time, greedy correctness proof.

\chapter{Fractional Knapsack}
Greedy by value-to-weight ratio, sorting step, O(n log n).

\chapter{Huffman Coding}
Optimal prefix-free codes, priority queue, tree construction from array.

\chapter{Backtracking Fundamentals}
Exhaustive search with pruning, state space tree, recursive exploration.

\chapter{N-Queens Problem}
Placing queens on chessboard, backtracking, array representation.

\chapter{Subset Generation}
Generating all subsets, backtracking, bit manipulation approaches.

\chapter{Permutation Generation}
Generating all permutations, backtracking, iterative algorithms (Heap's).

\chapter{Sudoku Solver}
Constraint satisfaction, backtracking, array representation.

% --- Online, streaming, and amortized analysis ---

\chapter{Online Algorithms: Introduction}
Online vs. offline, competitive ratio, worst-case analysis.

\chapter{Online Algorithms: Paging}
Cache replacement (LRU, FIFO, optimal), competitive analysis.

\chapter{Streaming Algorithms: Fundamentals}
One-pass algorithms, sublinear space, sampling techniques.

\chapter{Streaming: Heavy Hitters}
Finding frequent elements, Count-Min sketch, Misra-Gries algorithm.

\chapter{Streaming: Distinct Elements}
Flajolet-Martin, HyperLogLog, approximate counting.

\chapter{Reservoir Sampling}
Uniform random sample, single-pass, unknown stream size.

\chapter{Amortized Analysis: Aggregate Method}
Total cost divided by operations, examples (dynamic array, binary counter).

\chapter{Amortized Analysis: Accounting Method}
Amortized cost as credit, maintaining non-negative balance.

\chapter{Amortized Analysis: Potential Method}
Potential function, relating amortized to actual cost, flexibility.

\chapter{Cache-Oblivious Algorithms}
Optimal cache usage without knowing cache parameters, ideal cache model.

\chapter{Cache-Oblivious Matrix Transpose}
Recursive divide-and-conquer, automatically adapting to cache hierarchy.

\chapter{Cache-Oblivious Searching and Sorting}
Van Emde Boas layout, funnelsort, optimal cache complexity.

% ============================================================
% SECTION F: Parallel and distributed array processing
%   (merged from old Part 5, algorithm-level chapters only;
%    pure-hardware GPU/SIMD chapters already live in Volume II)
% ============================================================

\chapter{Concurrency vs. Parallelism}
Definitions, distinctions, concurrent execution, parallel execution, when each matters.

\chapter{Parallel Performance Metrics}
Speedup, efficiency, scalability, work, span, Amdahl's law, Gustafson's law.

\chapter{Amdahl's Law}
Serial fraction limits speedup, diminishing returns, implications for algorithm design.

\chapter{Gustafson's Law}
Scaled speedup, weak scaling vs. strong scaling, problem size growth.

\chapter{Work and Span Model}
DAG representation, work (total operations), span (critical path), parallelism = work/span.

\chapter{Processes vs. Threads}
Separate address spaces, shared memory, creation overhead, context switching.

\chapter{Thread Creation and Management}
POSIX threads (pthreads), C++11 std::thread, thread creation, joining, detaching.

\chapter{Thread Pools}
Worker threads, task queues, reducing creation overhead, load distribution.

\chapter{Work Stealing}
Dynamic load balancing, task stealing from other threads, Cilk model, deque-based implementation.

\chapter{Fork-Join Parallelism}
Divide and conquer parallelism, recursive task spawning, join synchronization.

\chapter{Parallel Recursion}
Parallelizing divide-and-conquer algorithms, base case threshold, overhead management.

\chapter{OpenMP: Parallel Loops}
Parallel for directive, work sharing, scheduling (static, dynamic, guided), loop dependencies.

\chapter{OpenMP: Data Sharing Clauses}
Shared, private, firstprivate, lastprivate, reduction, default clauses.

\chapter{OpenMP: Reductions}
Reduction operations, private copies, combining results, common patterns.

\chapter{OpenMP: Task Parallelism}
Task directive, task dependencies, taskwait, taskloop, irregular parallelism.

\chapter{OpenMP: SIMD Directives}
SIMD loops, aligned clauses, vectorization hints, safelen.

\chapter{OpenMP: Memory Model}
Flush operations, memory ordering, atomic operations, thread-safe programming.

\chapter{Data Races and Race Conditions}
Concurrent read-write, write-write conflicts, non-deterministic behavior, debugging races.

\chapter{Memory Consistency Models}
Sequential consistency, total store order (TSO), weak ordering, relaxed models.

\chapter{Happens-Before Relation}
Partial ordering, synchronizes-with, inter-thread ordering, reasoning about concurrency.

\chapter{Memory Ordering in C++11}
memory\_order\_relaxed, memory\_order\_acquire, memory\_order\_release, memory\_order\_seq\_cst.

\chapter{Synchronization Primitives: Mutexes}
Mutual exclusion locks, lock/unlock, critical sections, deadlock potential.

\chapter{Spinlocks vs. Mutexes}
Busy waiting vs. blocking, context switch overhead, when to use each.

\chapter{Reader-Writer Locks}
Multiple readers, exclusive writer, read-heavy workloads, fairness considerations.

\chapter{Semaphores}
Counting semaphores, binary semaphores, signaling, producer-consumer synchronization.

\chapter{Barriers}
Thread synchronization point, waiting for all threads, reusable barriers, OpenMP barrier.

\chapter{Condition Variables}
Wait/notify pattern, avoiding busy waiting, producer-consumer queues, spurious wakeups.

\chapter{Lock-Free Programming: Fundamentals}
Atomic operations, compare-and-swap (CAS), ABA problem, wait-free vs. lock-free.

\chapter{Atomic Operations}
Atomic reads/writes, fetch-and-add, compare-and-exchange, memory ordering.

\chapter{Lock-Free Stack}
CAS-based push/pop, ABA problem, tagged pointers, hazard pointers.

\chapter{Lock-Free Queue}
Michael-Scott queue, enqueue/dequeue with CAS, memory reclamation challenges.

\chapter{Memory Reclamation in Lock-Free Structures}
Hazard pointers, epoch-based reclamation, reference counting, safe memory reclamation (SMR).

\chapter{Transactional Memory: Concepts}
Atomic blocks, optimistic concurrency, conflict detection, rollback.

\chapter{Software Transactional Memory (STM)}
STM implementations, versioning, validation, commit protocol.

\chapter{Hardware Transactional Memory (HTM)}
Intel TSX, cache-based implementation, capacity limits, fallback paths.

\chapter{False Sharing}
Cache line contention, performance degradation, detecting false sharing (perf, VTune).

% --- GPU array algorithms (algorithm-level; GPU hardware lives in Volume II) ---

\chapter{GPU Sorting Algorithms}
Bitonic sort, radix sort on GPU, merge sort, library implementations (Thrust).

\chapter{GPU Reduction Algorithms}
Tree-based reduction, warp-level primitives, shuffle operations, optimized patterns.

\chapter{GPU Scan (Prefix Sum)}
Blelloch scan on GPU, work-efficient implementation, bank conflict avoidance.

\chapter{GPU Matrix Operations}
Matrix multiplication, tiling, shared memory usage, cuBLAS library.

\chapter{GPU Stencil Computations}
Halo exchange, shared memory tiling, minimizing global memory access.

\chapter{Heterogeneous Computing}
CPU-GPU collaboration, task partitioning, host-device coordination.

\chapter{OpenCL Programming Model}
Platform, device, context, command queue, kernels, portability across devices.

\chapter{OpenCL Memory Model}
Global, local, private, constant memory, explicit transfers, buffer objects.

% --- Distributed array processing ---

\chapter{Distributed Computing Fundamentals}
Distributed memory, message passing, latency and bandwidth, network topology.

\chapter{Message Passing Interface (MPI): Basics}
Communicators, ranks, MPI\_Init, MPI\_Finalize, basic communication.

\chapter{MPI Point-to-Point Communication}
Send, receive, buffering modes, blocking vs. non-blocking, message tags.

\chapter{MPI Non-Blocking Communication}
MPI\_Isend, MPI\_Irecv, overlapping communication and computation, MPI\_Wait.

\chapter{MPI Collective Operations}
Broadcast, scatter, gather, all-to-all, reduction, barrier, synchronization.

\chapter{MPI Reductions}
MPI\_Reduce, MPI\_Allreduce, custom reduction operations, combining data.

\chapter{MPI Derived Datatypes}
Contiguous, vector, indexed types, describing complex data layouts.

\chapter{Distributed Array Partitioning}
Block distribution, cyclic distribution, block-cyclic, load balancing.

\chapter{Ghost Cells and Halo Exchange}
Boundary regions, stencil computations, communication patterns, optimizing exchange.

\chapter{Domain Decomposition}
1D, 2D, 3D decomposition, minimizing surface-to-volume ratio, communication overhead.

\chapter{Load Balancing in Distributed Systems}
Static vs. dynamic, work stealing across nodes, balancing irregular workloads.

\chapter{Parallel I/O}
MPI-IO, parallel file systems (Lustre, GPFS), collective I/O, avoiding bottlenecks.

\chapter{Checkpoint/Restart}
Fault tolerance, periodic checkpointing, restart from checkpoint, overhead.

\chapter{Fault Tolerance Strategies}
Replication, algorithm-based fault tolerance, detecting failures, recovery.

\chapter{MapReduce Programming Model}
Map phase, shuffle/sort, reduce phase, data-parallel programming at scale.

\chapter{MapReduce on Arrays}
Array processing patterns, sorting, filtering, aggregation, distributed arrays.

\chapter{Hadoop and HDFS}
Distributed file system, block storage, replication, data locality.

\chapter{Spark and RDDs}
Resilient distributed datasets, in-memory processing, transformations and actions.

\chapter{Spark Array Operations}
DataFrame operations, array columns, broadcast variables, distributed array processing.

\chapter{Partitioned Global Address Space (PGAS)}
Global view with local control, UPC, Chapel, X10, one-sided communication.

\chapter{Array Notation in PGAS Languages}
Distributed array declarations, affinity, local vs. remote access.

\chapter{Distributed Consensus}
Paxos, Raft, leader election, maintaining consistency, fault tolerance.

\chapter{Eventual Consistency}
Convergent data types, CRDTs, conflict resolution, distributed arrays.

% --- Real-time and energy-aware array processing ---

\chapter{Real-Time Systems}
Hard vs. soft deadlines, WCET (worst-case execution time), predictability.

\chapter{WCET Analysis for Array Operations}
Static analysis, timing bounds, cache analysis, array access time.

\chapter{Real-Time Scheduling}
Rate monotonic scheduling (RMS), earliest deadline first (EDF), priority assignment.

\chapter{Cache Locking for Real-Time}
Locking cache lines, predictable access times, avoiding cache misses.

\chapter{Scratchpad Memory}
Software-managed memory, predictable access, explicitly controlled, vs. caches.

\chapter{Energy-Efficient Computing}
Power models, dynamic power, static power, energy-aware algorithms.

\chapter{DVFS and Power Management}
Dynamic voltage and frequency scaling, race-to-idle vs. energy proportional.

\chapter{Energy-Efficient Array Processing}
Minimizing memory access, data reuse, approximate computing for energy savings.

\chapter{Green Computing Principles}
Energy proportionality, workload consolidation, resource efficiency.

\chapter{Parallel Debugging}
Race detection, Valgrind, ThreadSanitizer, Intel Inspector, deterministic replay.

\chapter{Performance Analysis Tools}
gprof, perf, VTune, TAU, HPCToolkit, profiling parallel programs.

\chapter{Scalability Analysis}
Strong scaling, weak scaling, scalability plots, identifying bottlenecks.

\chapter{Communication vs. Computation Ratio}
Bandwidth-limited vs. compute-limited, optimizing the ratio, hiding latency.

\chapter{Parallel Algorithm Design Patterns}
Task parallelism, data parallelism, pipeline, master-worker, SPMD.

% ============================================================
% SECTION G: Arrays in science, machine learning, and beyond
%   (merged from old Part 6)
% ============================================================

\chapter{Tensors and Multidimensional Arrays}
Mathematical tensors, rank/order, tensor operations, Einstein summation notation.

\chapter{Tensor Computing Frameworks}
NumPy, TensorFlow, PyTorch, JAX, computational graphs, automatic differentiation.

\chapter{Broadcasting Semantics}
NumPy broadcasting rules, implicit dimension expansion, shape compatibility.

\chapter{Einstein Summation (einsum)}
Concise tensor operations, einsum notation, optimization, implementing einsum.

\chapter{Neural Networks as Array Transformations}
Weight matrices, bias vectors, activation functions, layer operations.

\chapter{Convolutional Neural Networks}
Convolution operation, filters as arrays, feature maps, receptive fields.

\chapter{Convolution Implementation}
Direct convolution, im2col transformation, Winograd convolution, FFT-based convolution.

\chapter{Pooling Operations}
Max pooling, average pooling, global pooling, downsampling arrays.

\chapter{Batch Normalization}
Normalizing activations, running statistics, batch dimension, inference mode.

\chapter{Dropout and Regularization}
Random array masking, training vs. inference, regularization effects.

\chapter{Backpropagation as Array Operations}
Chain rule, gradient computation, forward and backward passes, computational graph.

\chapter{Automatic Differentiation}
Forward mode, reverse mode (backpropagation), dual numbers, computational graph.

\chapter{Gradient Arrays and Optimizers}
SGD, momentum, Adam, RMSprop, gradient clipping, weight updates as array operations.

\chapter{Mini-Batch Training}
Batching samples, parallel processing, batch size effects, memory considerations.

\chapter{Data Loading Pipelines}
Prefetching, data augmentation, batching strategies, avoiding bottlenecks.

\chapter{Model Parallelism}
Splitting models across devices, pipeline parallelism, tensor parallelism.

\chapter{Data Parallelism in Training}
Replicating models, gradient synchronization, all-reduce, parameter servers.

\chapter{Quantization Techniques}
Post-training quantization, quantization-aware training, INT8, mixed precision.

\chapter{Neural Network Pruning}
Weight pruning, structured vs. unstructured, sparse neural networks, compression.

\chapter{Knowledge Distillation}
Teacher-student models, soft targets, model compression, array operations.

\chapter{Attention Mechanisms}
Query-key-value arrays, scaled dot-product attention, multi-head attention.

\chapter{Transformer Architecture}
Self-attention, positional encoding, encoder-decoder, array shapes in transformers.

\chapter{Sparse Attention Patterns}
Reducing quadratic complexity, local attention, strided attention, efficient transformers.

\chapter{Tensor Cores and Hardware Acceleration}
Mixed-precision matrix multiplication, Tensor Core operations (NVIDIA), TPU systolic arrays.

% --- Linear algebra and scientific computing ---

\chapter{Matrix Multiplication Optimization}
Blocking/tiling, Strassen's algorithm, cache optimization, BLAS libraries.

\chapter{BLAS and LAPACK}
Basic Linear Algebra Subprograms, Level 1/2/3 BLAS, LAPACK routines, vendor implementations.

\chapter{Dense Linear Algebra}
Matrix decompositions (LU, QR, Cholesky, SVD), eigenvalue problems, numerical stability.

\chapter{Sparse Linear Algebra}
Sparse matrix-vector multiply (SpMV), sparse solvers, iterative methods, preconditioning.

\chapter{Iterative Solvers}
Jacobi, Gauss-Seidel, conjugate gradient, GMRES, array-based iterations.

\chapter{Scientific Computing: Finite Differences}
Discretizing derivatives, stencil operations, explicit vs. implicit methods, stability.

\chapter{Finite Element Methods}
Mesh representation, element arrays, assembly, sparse matrix systems.

\chapter{Computational Fluid Dynamics}
Grid-based simulations, array representations, parallel CFD, large-scale simulations.

\chapter{N-Body Simulations}
Particle arrays, force calculations, Barnes-Hut algorithm, fast multipole method.

% --- Signal and media processing ---

\chapter{Fast Fourier Transform (FFT)}
Cooley-Tukey algorithm, radix-2, DIT vs. DIF, bit-reversal permutation.

\chapter{FFT Applications}
Signal processing, convolution via FFT, polynomial multiplication, audio processing.

\chapter{Multidimensional FFT}
2D and 3D FFT, separable transforms, image processing, spectral methods.

\chapter{Image Processing Fundamentals}
Images as 2D/3D arrays, pixels, color channels, spatial operations.

\chapter{Image Filtering}
Convolution with kernels, Gaussian blur, edge detection (Sobel, Canny), median filter.

\chapter{Morphological Operations}
Erosion, dilation, opening, closing, structuring elements, binary image processing.

\chapter{Image Transformations}
Rotation, scaling, translation, affine transformations, interpolation methods.

\chapter{Feature Detection}
Corner detection (Harris), SIFT, SURF, ORB, keypoint arrays.

\chapter{Video Processing}
Frame arrays, temporal dimension, motion estimation, optical flow.

\chapter{Audio Signal Processing}
Waveform arrays, sampling, windowing, STFT (short-time Fourier transform).

\chapter{Digital Filters}
FIR filters, IIR filters, filter design, convolution implementation, Z-transform.

\chapter{Spectral Analysis}
Power spectral density, spectrogram, frequency domain analysis, windowing effects.

\chapter{Time Series Analysis}
Temporal arrays, trend analysis, seasonal decomposition, ARIMA models.

\chapter{Time Series Forecasting}
Moving averages, exponential smoothing, LSTM networks, array-based forecasting.

\chapter{Wavelet Transforms}
Continuous vs. discrete wavelets, multiresolution analysis, image compression.

% --- Geospatial and bioinformatics ---

\chapter{Geospatial Data Arrays}
Raster data, elevation maps, GIS operations, spatial indexing.

\chapter{Remote Sensing}
Satellite imagery, multispectral arrays, band operations, image classification.

\chapter{Climate and Weather Models}
Grid-based simulations, atmospheric models, large-scale array computations.

\chapter{Genomic Sequence Arrays}
DNA sequences as character arrays, k-mers, sequence alignment algorithms.

\chapter{Genome Assembly}
De Bruijn graphs, overlap-layout-consensus, array-based algorithms, parallelization.

\chapter{Sequence Alignment}
Smith-Waterman, BLAST, dynamic programming on arrays, scoring matrices.

\chapter{Gene Expression Analysis}
Expression matrices, differential expression, clustering, heatmaps.

\chapter{Proteomics and Mass Spectrometry}
Spectral arrays, peak detection, database searching, peptide identification.

\chapter{Molecular Dynamics Simulations}
Particle position/velocity arrays, force calculations, integration methods.

% --- Quantum computing ---

\chapter{Quantum Computing Fundamentals}
Qubits, quantum gates, quantum circuits, measurement, superposition.

\chapter{Quantum State Vectors}
Complex-valued arrays, state representation, probability amplitudes, normalization.

\chapter{Quantum Gates as Matrix Operations}
Pauli gates, Hadamard, CNOT, unitary matrices, gate decomposition.

\chapter{Quantum Algorithms}
Grover's search, Shor's factoring, quantum Fourier transform, amplitude amplification.

\chapter{Quantum Circuit Simulation}
Statevector simulation, density matrices, simulating quantum computers classically.

% --- DNA storage and emerging paradigms ---

\chapter{DNA Data Storage}
Encoding binary data in DNA sequences, error correction, reading/writing DNA.

\chapter{DNA Storage: Array Representation}
Mapping bits to nucleotides, redundancy, addressing, random access.

\chapter{Neuromorphic Computing}
Spiking neural networks, event-driven computation, neuromorphic chips (Loihi, TrueNorth).

\chapter{Spike Arrays}
Temporal spike patterns, spike-timing-dependent plasticity, encoding information.

\chapter{Analog Computing Arrays}
Memristor crossbar arrays, analog matrix multiplication, in-memory computing.

\chapter{Processing-in-Memory (PIM)}
Reducing data movement, near-memory computation, avoiding von Neumann bottleneck.

% --- Array programming languages and ecosystems ---

\chapter{Array Programming Languages}
APL, J, K, Q, array-oriented paradigm, implicit parallelism, concise notation.

\chapter{APL and J Notation}
Array operations, rank polymorphism, operators vs. functions, reading APL.

\chapter{Julia for Array Computing}
Multiple dispatch, broadcasting, SIMD, GPU support, scientific computing.

\chapter{NumPy Internals}
ndarray structure, strides, memory layout, vectorized operations, ufuncs.

\chapter{NumPy Performance Optimization}
Vectorization, avoiding loops, broadcasting, numba JIT compilation.

\chapter{Pandas and DataFrames}
Columnar data, array backing, operations on arrays, time series support.

\chapter{Dask for Parallel Arrays}
Lazy evaluation, parallel and out-of-core arrays, distributed computing, task graphs.

\chapter{Array Databases}
SciDB, RasdaMan, TileDB, array query languages, chunking, distribution.

\chapter{Column-Store Databases}
Columnar storage, vectorized query execution, compression, analytical workloads.

\chapter{Vectorized Query Processing}
Batch processing, avoiding row-at-a-time, SIMD in databases, MonetDB.

\chapter{Just-In-Time Compilation for Arrays}
Numba, JAX jit, PyPy, runtime optimization, trace compilation.

\chapter{XLA Compiler}
Accelerated linear algebra, TensorFlow/JAX backend, fusion, optimization passes.

\chapter{Polyhedral Optimization}
Loop transformations, dependence analysis, automatic parallelization, tiling.

% --- Out-of-core, persistent, and compressed array representations ---

\chapter{Memory-Mapped Arrays}
mmap, virtual memory tricks, out-of-core computation, file-backed arrays.

\chapter{Out-of-Core Array Processing}
Working with arrays larger than RAM, streaming algorithms, blocked processing.

\chapter{Persistent Data Structures}
Immutable arrays, structural sharing, path copying, functional programming.

\chapter{Versioned Arrays}
Copy-on-write, tracking versions, git-like array versioning, efficient diffs.

\chapter{Compressed Arrays: Run-Length Encoding}
Consecutive duplicates, compression ratio, when effective, sparse data.

\chapter{Compressed Arrays: Dictionary Encoding}
Categorical data, integer codes, lookup table, columnar databases.

\chapter{Compressed Arrays: Delta Encoding}
Storing differences, time series, sorted data, integer compression.

\chapter{Compressed Arrays: Frame-of-Reference}
Base value plus small deltas, bit packing, database compression.

\chapter{Succinct Data Structures}
Near-optimal space, rank and select queries, Jacobson's rank structure.

\chapter{Wavelet Trees}
Multiary wavelet tree, range queries, succinct representation, text indexing.

\chapter{Cache-Oblivious Algorithms Revisited}
Recursive algorithms, automatically optimal, no tuning, ideal cache model.

\chapter{Van Emde Boas Layout}
Recursive array layout, cache-oblivious B-trees, implicit data structures.

\chapter{External Memory Algorithms}
I/O model, sorting, searching, graph algorithms, minimizing block transfers.

\chapter{Streaming Algorithms}
Sublinear space, one pass, sampling, sketching, space-time tradeoffs.

\chapter{Approximate Computing}
Trading accuracy for performance/energy, probabilistic guarantees, acceptable error bounds.

\chapter{Probabilistic Data Structures Summary}
Bloom filters, Count-Min sketch, HyperLogLog, comparison of techniques.

% --- Privacy, security, and distributed-ledger array structures ---

\chapter{Privacy-Preserving Array Operations}
Homomorphic encryption, secure multi-party computation, differential privacy.

\chapter{Differential Privacy}
Adding noise, privacy budget, Laplace mechanism, Gaussian mechanism, query responses.

\chapter{Homomorphic Encryption on Arrays}
Computing on encrypted data, fully homomorphic encryption (FHE), performance challenges.

\chapter{Secure Multi-Party Computation}
Secret sharing, garbled circuits, oblivious transfer, distributed array processing.

\chapter{Blockchain and Merkle Trees}
Hash trees over arrays, integrity proofs, efficient verification, cryptocurrency applications.

\chapter{Distributed Ledger Arrays}
Append-only arrays, consensus, Byzantine fault tolerance, state replication.

% --- Edge, embedded, and emerging hardware paradigms ---

\chapter{Edge Computing and IoT}
Resource-constrained devices, local array processing, edge inference, federated learning.

\chapter{Federated Learning}
Distributed training, privacy preservation, aggregating gradients, edge devices.

\chapter{TinyML and Embedded Arrays}
Running neural networks on microcontrollers, quantization, model compression, real-time inference.

\chapter{Hardware Accelerators for Arrays}
FPGAs, ASICs, domain-specific architectures, spatial architectures.

\chapter{FPGA Array Processing}
Reconfigurable logic, custom datapaths, pipelining, HLS (high-level synthesis).

\chapter{Optical Computing}
Photonic computing, optical matrix multiplication, speed-of-light computation.

\chapter{Reversible Computing}
Landauer limit, reversible logic, energy efficiency, adiabatic computing.

\chapter{Biological Computing}
DNA computing, cellular automata, biological circuits, array-like structures.

\chapter{Future Memory Technologies}
Persistent memory (3D XPoint), NVRAM, storage-class memory, implications for arrays.

\chapter{Post-Moore's Law Computing}
3D stacking, chiplets, advanced packaging, specialized architectures, domain-specific computing.

\chapter{Array Computing: Future Directions}
Research frontiers, open problems, emerging paradigms, the next 50 years.